\chapter{Green, Stokes and Gauss}
\label{ch:theorems}

This chapter contains the three great integral theorems of vector calculus. Each one says the same thing in a different dimension:
\[
\text{(integral of a derivative over a region)}=\text{(integral of the function over the boundary of the region)} .
\]
This is the fundamental theorem of calculus, $\int_a^bf'\,\dd x=f(b)-f(a)$, in higher dimensions. We prove each theorem for a basic class of regions, and then explain how the general case follows by cutting a region into basic pieces.

%====================================================================
\section{Green's theorem in the plane}
%====================================================================

Let $D\subset\R^2$ be a bounded region whose boundary $\partial D$ is a piecewise $C^1$ closed curve (or several of them). The boundary is given the \emph{positive orientation}: when you walk along it, the region $D$ is on your \emph{left}. For a disc this means counterclockwise; for a region with a hole, the outer boundary runs counterclockwise and the inner boundary clockwise.

\begin{theorem}[Green's theorem]\label{thm:green}
Let $P,Q$ be $C^1$ functions on an open set containing $D\cup\partial D$. Then
\[
\oint_{\partial D}P\,\dd x+Q\,\dd y=\iint_D\Bigl(\pd Qx-\pd Py\Bigr)\,\dd A .
\]
\end{theorem}

In vector form, with $\vec F=(P,Q)$: $\oint_{\partial D}\vec F\cdot\dd\vec r=\iint_D(\curl\vec F)\cdot\khat\,\dd A$. The circulation around the boundary equals the total of the local circulation densities (curl) inside.

\begin{figure}[ht]
\centering
\begin{tikzpicture}[scale=1.1]
  \begin{scope}
    \filldraw[fill=accent!12,draw=accent,thick,postaction={decorate,decoration={markings,mark=at position 0.3 with {\arrow{Stealth}},mark=at position 0.8 with {\arrow{Stealth}}}}]
      (0,0) to[out=-20,in=200] (3.2,0.2) to[out=60,in=-100] (3.5,2.4) to[out=120,in=20] (0.3,2.6) to[out=200,in=100] cycle;
    \node at (1.7,1.3) {$D$};
    \node[accent,font=\small] at (3.9,1.2) {$\partial D$};
  \end{scope}
  \begin{scope}[xshift=5.6cm]
    \filldraw[fill=accent!12,draw=accent,thick] (0.2,0) rectangle (3.4,2.6);
    \filldraw[fill=white,draw=accent2,thick] (1.3,0.8) rectangle (2.3,1.8);
    \draw[accent,very thick,->] (0.2,0)--(0.2,0.01);
    \draw[accent,thick,-{Stealth}] (1.5,0)--(2.1,0);
    \draw[accent,thick,-{Stealth}] (3.4,1.2)--(3.4,1.5);
    \draw[accent,thick,-{Stealth}] (2.1,2.6)--(1.5,2.6);
    \draw[accent,thick,-{Stealth}] (0.2,1.5)--(0.2,1.2);
    \draw[accent2,thick,-{Stealth}] (1.3,1.5)--(1.3,1.1);
    \draw[accent2,thick,-{Stealth}] (1.5,1.8)--(1.9,1.8);
    \draw[accent2,thick,-{Stealth}] (2.3,1.1)--(2.3,1.5);
    \draw[accent2,thick,-{Stealth}] (2.1,0.8)--(1.7,0.8);
    \node at (2.8,2.0) {$D$};
    \node[font=\small,accent2] at (1.8,1.3) {hole};
  \end{scope}
\end{tikzpicture}
\caption{Positive orientation: the region is on the left as we walk. Left: a simple region, counterclockwise. Right: a region with a hole; the outer boundary is counterclockwise and the inner boundary is clockwise.}
\label{fig:greenorient}
\end{figure}

\begin{proof}[Proof for regions that are of both type I and type II]
\emph{Part 1: $\oint P\,\dd x=-\iint P_y\,\dd A$.} Write $D=\{(x,y):a\le x\le b,\ g_1(x)\le y\le g_2(x)\}$ with $g_1\le g_2$ continuous and piecewise $C^1$. By the fundamental theorem of calculus in the variable $y$ and Fubini (Theorem~\ref{thm:fubini}),
\[
\iint_D\pd Py\,\dd A=\int_a^b\bigl[P(x,g_2(x))-P(x,g_1(x))\bigr]\dd x .
\]
Now compute $\oint_{\partial D}P\,\dd x$. The boundary has four parts. The bottom curve $y=g_1(x)$ is traversed from $x=a$ to $x=b$ (left to right, with $D$ above, on the left): this contributes $\int_a^bP(x,g_1(x))\dd x$. The top curve $y=g_2(x)$ is traversed from $b$ to $a$: it contributes $-\int_a^bP(x,g_2(x))\dd x$. The two vertical pieces (where $x$ is constant) have $\dd x=0$ and contribute nothing. Hence
\[
\oint_{\partial D}P\,\dd x=\int_a^bP(x,g_1)\dd x-\int_a^bP(x,g_2)\dd x=-\iint_D\pd Py\,\dd A .
\]
\emph{Part 2: $\oint Q\,\dd y=\iint Q_x\,\dd A$.} Write $D=\{c\le y\le d,\ h_1(y)\le x\le h_2(y)\}$. Then $\iint_DQ_x\dd A=\int_c^d[Q(h_2(y),y)-Q(h_1(y),y)]\dd y$. The right curve $x=h_2(y)$ is traversed upward (from $c$ to $d$) and the left curve $x=h_1(y)$ downward, and the horizontal pieces have $\dd y=0$; so $\oint Q\,\dd y=\int_c^dQ(h_2,y)\dd y-\int_c^dQ(h_1,y)\dd y$, the same quantity.

Adding Parts 1 and 2 gives the theorem.
\end{proof}

\paragraph{From simple regions to general regions.} Suppose $D$ is cut by a curve into two regions $D_1,D_2$ for which Green's theorem holds. The cut is a part of the boundary of both pieces, but traversed in \emph{opposite directions} (each piece has itself on its left). So the line integrals along the cut cancel in the sum $\oint_{\partial D_1}+\oint_{\partial D_2}=\oint_{\partial D}$, and the double integrals add to $\iint_D$. Therefore Green's theorem holds for $D$. Repeating this, it holds for every region that can be cut into finitely many regions of both types: discs, annuli, polygons, regions with finitely many holes, and so on. The statement for \emph{arbitrary} piecewise smooth boundaries follows by approximation; we take it for granted.

\begin{figure}[ht]
\centering
\begin{tikzpicture}[scale=1.2]
  \filldraw[fill=accent!10,draw=none] (0,0) rectangle (4,2.4);
  \draw[accent,thick] (0,0) rectangle (4,2.4);
  \draw[accent2,very thick] (2,0)--(2,2.4);
  \draw[accent2,thick,-{Stealth}] (2.06,1.3)--(2.06,0.7);
  \draw[accent3,thick,-{Stealth}] (1.94,0.9)--(1.94,1.5);
  \node at (1,1.2) {$D_1$}; \node at (3,1.2) {$D_2$};
  \draw[accent,thick,-{Stealth}] (0.4,0)--(1.0,0); \draw[accent,thick,-{Stealth}] (2.4,0)--(3.0,0);
  \draw[accent,thick,-{Stealth}] (4,0.8)--(4,1.4);
  \draw[accent,thick,-{Stealth}] (3.0,2.4)--(2.4,2.4); \draw[accent,thick,-{Stealth}] (1.0,2.4)--(0.4,2.4);
  \draw[accent,thick,-{Stealth}] (0,1.4)--(0,0.8);
\end{tikzpicture}
\caption{The cut between $D_1$ and $D_2$ is traversed once in each direction. The two line integrals along it cancel, so the sum of the boundary integrals of the pieces is the boundary integral of the whole.}
\label{fig:cutcancel}
\end{figure}

\begin{example}[Area as a line integral]
Take $P=-y/2$, $Q=x/2$. Then $Q_x-P_y=1$, and Green's theorem gives
\[
\Area(D)=\frac12\oint_{\partial D}(x\,\dd y-y\,\dd x).
\]
For the ellipse $x=a\cos t$, $y=b\sin t$: $x\dd y-y\dd x=ab\,\dd t$, so $\Area=\frac12\cdot2\pi ab=\pi ab$. (This is the planimeter principle: running a wheel along the boundary measures the area.)
\end{example}

\begin{example}[Computing a hard line integral]
Let $C$ be the triangle with vertices $(0,0),(1,0),(0,1)$, counterclockwise, and $\vec F=(y^2,\ 3xy)$. Then $Q_x-P_y=3y-2y=y$ and
\[
\oint_C\vec F\cdot\dd\vec r=\int_0^1\!\!\int_0^{1-x}y\,\dd y\,\dd x=\int_0^1\frac{(1-x)^2}2\dd x=\frac16 .
\]
\end{example}

\subsection{The vortex revisited}

Green's theorem does not apply to the vortex field $\vec G$ of Chapter~\ref{ch:linesurface} on a region containing the origin, because $\vec G$ is not $C^1$ there. But it does apply to an annulus $D=\{\varepsilon\le r\le1\}$, where $\curl\vec G=0$. The boundary consists of the unit circle (counterclockwise) and the small circle (clockwise), so
\[
0=\oint_{r=1}\vec G\cdot\dd\vec r-\oint_{r=\varepsilon}\vec G\cdot\dd\vec r,
\]
where both circles are now taken counterclockwise. The circulation is the same ($2\pi$) on every circle around the origin. This is the idea of ``deforming the contour'' that is used all over mathematics and physics.

\begin{theorem}[Green's theorem for a field with a point singularity]\label{thm:greensing}
Let $\vec F$ be $C^1$ with $\curl\vec F\cdot\khat=0$ on $D\setminus\{\vec p\}$, where $\vec p$ is an interior point of $D$. Then $\oint_{\partial D}\vec F\cdot\dd\vec r=\oint_{C_\varepsilon}\vec F\cdot\dd\vec r$ where $C_\varepsilon$ is any small counterclockwise circle around $\vec p$.
\end{theorem}

\begin{proof}
Apply Green to $D\setminus B_\varepsilon(\vec p)$, which has boundary $\partial D$ together with $C_\varepsilon$ reversed.
\end{proof}

%====================================================================
\section{The divergence theorem (Gauss)}
%====================================================================

Let $E\subset\R^3$ be a bounded solid region whose boundary $\partial E$ is a closed piecewise smooth surface. The \emph{outward} unit normal is denoted $\hat{\vec n}$.

\begin{theorem}[Divergence theorem]\label{thm:gauss}
Let $\vec F=(P,Q,R)$ be $C^1$ on an open set containing $E\cup\partial E$. Then
\[
\iiint_E\divg\vec F\,\dd V=\oiint_{\partial E}\vec F\cdot\hat{\vec n}\,\dd S .
\]
\end{theorem}

(We write $\oiint$ for the surface integral over a closed surface; here it is only a notation for $\iint$ with a circle.) The total source strength inside equals the net flux out through the boundary.

\begin{figure}[ht]
\centering
\begin{tikzpicture}[scale=1.15]
  \path[fill=accent!10] (0,0.3) to[out=60,in=180] (1.6,2.3) to[out=0,in=140] (3.6,1.9) to[out=-40,in=70] (4.2,0.6) to[out=-110,in=10] (2.6,-0.3) to[out=190,in=-60] cycle;
  \draw[accent,thick] (0,0.3) to[out=60,in=180] (1.6,2.3) to[out=0,in=140] (3.6,1.9) to[out=-40,in=70] (4.2,0.6) to[out=-110,in=10] (2.6,-0.3) to[out=190,in=-60] cycle;
  \node at (2.1,1.0) {$E$};
  \node[accent,font=\small] at (3.4,2.4) {$\partial E$};
  \foreach \p/\a in {{0.55,1.35}/150,{1.6,2.28}/90,{3.0,2.22}/45,{4.1,1.3}/10,{3.9,0.1}/-60,{1.5,-0.1}/-100}{
    \draw[vecstyle,accent2] (\p)--++(\a:0.55);
  }
  \node[accent2,font=\small] at (4.9,1.55) {$\hat{\vec n}$};
  % small sources inside
  \foreach \p in {(1.4,0.9),(2.1,1.6),(2.9,0.7)}{ \node[accent3] at \p {$+$}; }
\end{tikzpicture}
\caption{The divergence theorem: the net outward flux through the closed surface is the total of the sources (divergence) inside.}
\label{fig:gauss}
\end{figure}

\begin{proof}[Proof for a region simple in every direction]
Write $\vec F=P\ihat+Q\jhat+R\khat$. Since $\divg\vec F=P_x+Q_y+R_z$ and the surface integral is $\oiint(P\,n_x+Q\,n_y+R\,n_z)\dd S$, it suffices to prove the three identities
\[
\iiint_EP_x\,\dd V=\oiint P\,n_x\,\dd S,\qquad\iiint_EQ_y\dd V=\oiint Q\,n_y\dd S,\qquad\iiint_ER_z\dd V=\oiint R\,n_z\dd S .
\]
We prove the third; the others are identical with the roles of the axes changed. Suppose
\[
E=\{(x,y,z):(x,y)\in D,\ g_1(x,y)\le z\le g_2(x,y)\}
\]
where $D$ is the projection of $E$ to the $xy$-plane and $g_1\le g_2$ are $C^1$ (a region \emph{simple in the $z$-direction}). By the fundamental theorem of calculus,
\[
\iiint_ER_z\,\dd V=\iint_D\bigl[R(x,y,g_2(x,y))-R(x,y,g_1(x,y))\bigr]\dd A .
\]
Now compute the surface integral. The boundary has three parts.
\begin{itemize}
  \item \emph{Top} $S_2:z=g_2(x,y)$, outward normal pointing up. By Theorem~\ref{thm:fluxgraph}, $\iint_{S_2}R\,n_z\,\dd S=\iint_DR(x,y,g_2)\,\dd A$ (take $\vec F=(0,0,R)$: the flux through an upward-oriented graph is $\iint R\,\dd A$).
  \item \emph{Bottom} $S_1:z=g_1(x,y)$, outward normal pointing down, i.e.\ opposite to the upward orientation: $\iint_{S_1}R\,n_z\,\dd S=-\iint_DR(x,y,g_1)\,\dd A$.
  \item \emph{Side} (vertical wall above $\partial D$, if any): the normal is horizontal, so $n_z=0$ and the integral is $0$.
\end{itemize}
The sum of the three is exactly the double integral above.
\end{proof}

A general region (a ball, a cube, a torus, a solid with a hole) is cut into finitely many regions that are simple in all three directions. The fluxes through the cut surfaces cancel, because the outward normals of the two neighbouring pieces are opposite. This is the same cancellation argument as for Green's theorem.

\begin{example}[Flux through a closed surface]
Compute the outward flux of $\vec F=(x^3,y^3,z^3)$ through the unit sphere. By the theorem, it equals $\iiint3(x^2+y^2+z^2)\dd V=3\int_0^1\rho^2\cdot4\pi\rho^2\,\dd\rho=\frac{12\pi}5$. A direct computation would be much longer.
\end{example}

\begin{example}[Volume from a surface integral]
For $\vec F=\frac13\vec r$, $\divg\vec F=1$ and so $\Vol(E)=\frac13\oiint_{\partial E}\vec r\cdot\hat{\vec n}\,\dd S$. For a ball of radius $R$, $\vec r\cdot\hat{\vec n}=R$ and $\Vol=\frac13R\cdot4\pi R^2=\frac43\pi R^3$.
\end{example}

\begin{example}[Gauss's law for a point charge]
Let $\vec F=\vec r/r^3$, which has $\divg\vec F=0$ for $\vec r\ne\vec0$. Let $E$ be the region between a small sphere $S_\varepsilon$ and a large closed surface $S$ around the origin. Then $0=\iiint_E\divg\vec F\,\dd V=\oiint_S\vec F\cdot\hat{\vec n}\dd S-\oiint_{S_\varepsilon}\vec F\cdot\hat{\vec n}\,\dd S$, where both normals are now pointing away from the origin. Since the flux through $S_\varepsilon$ is $4\pi$ (Chapter~\ref{ch:linesurface}), we get
\[
\oiint_S\frac{\vec r}{r^3}\cdot\hat{\vec n}\,\dd S=4\pi\qquad\text{for every closed surface $S$ enclosing the origin.}
\]
(And $0$ if the origin lies outside $S$.) This is the mathematical content of Gauss's law.
\end{example}

\begin{theorem}[Divergence as flux density, again]
For a $C^1$ field $\vec F$ and a point $\vec p$, $\divg\vec F(\vec p)=\lim_{\varepsilon\to0}\dfrac1{\Vol(B_\varepsilon)}\oiint_{\partial B_\varepsilon(\vec p)}\vec F\cdot\hat{\vec n}\,\dd S$.
\end{theorem}

\begin{proof}
By the divergence theorem, the right side is the average of $\divg\vec F$ over the ball $B_\varepsilon(\vec p)$, which tends to $\divg\vec F(\vec p)$ by continuity. (In Chapter~\ref{ch:fields} we proved this using a small box, \emph{before} we had the theorem; here it is a one-line consequence.)
\end{proof}

%====================================================================
\section{Stokes' theorem}
%====================================================================

Let $S$ be an oriented piecewise smooth surface in $\R^3$ with unit normal $\hat{\vec n}$, and let $\partial S$ be its boundary curve. The orientation of $\partial S$ is determined by the \emph{right-hand rule}: if you curl the fingers of your right hand in the direction of travel along $\partial S$, your thumb points along $\hat{\vec n}$. Equivalently, walking along $\partial S$ with your head pointing in the direction of $\hat{\vec n}$, the surface is on your left.

\begin{theorem}[Stokes' theorem]\label{thm:stokes}
Let $\vec F$ be $C^1$ on an open set containing $S$. Then
\[
\oint_{\partial S}\vec F\cdot\dd\vec r=\iint_S(\curl\vec F)\cdot\hat{\vec n}\,\dd S .
\]
\end{theorem}

\begin{figure}[ht]
\centering
\begin{tikzpicture}[scale=1.2]
  \path[fill=accent!12] (0,0.4) to[out=-30,in=200] (3.0,0.3) to[out=20,in=-120] (4.2,1.6) to[out=100,in=0] (3.0,2.6) to[out=180,in=60] (0.6,2.2) to[out=-130,in=60] cycle;
  \draw[accent,thick,postaction={decorate,decoration={markings,mark=at position 0.25 with {\arrow{Stealth}},mark=at position 0.7 with {\arrow{Stealth}}}}]
    (0,0.4) to[out=-30,in=200] (3.0,0.3) to[out=20,in=-120] (4.2,1.6) to[out=100,in=0] (3.0,2.6) to[out=180,in=60] (0.6,2.2) to[out=-130,in=60] cycle;
  \node at (2.2,1.4) {$S$};
  \node[accent,font=\small] at (4.5,0.5) {$\partial S$};
  \draw[vecstyle,accent2] (1.4,1.2)--++(0,1.1) node[above,font=\small]{$\hat{\vec n}$};
  \draw[accent4,-{Stealth},thick] (1.55,1.0) arc[start angle=0,end angle=300,x radius=0.0,y radius=0.0];
  \draw[accent4,thick,-{Stealth}] (1.1,1.35) arc[start angle=200,end angle=520,x radius=0.35,y radius=0.17];
\end{tikzpicture}
\caption{Stokes' theorem: the circulation around the edge $\partial S$ equals the flux of the curl through the surface. The direction on $\partial S$ is fixed by the right-hand rule relative to $\hat{\vec n}$.}
\label{fig:stokes}
\end{figure}

\begin{proof}[Proof for a surface that is the graph of a $C^2$ function]
Let $S$ be the graph $z=f(x,y)$ over a region $D$ for which Green's theorem holds, oriented with upward normal. Then $\partial S$ projects onto $\partial D$ with the same (counterclockwise) orientation. Write $\vec F=(P,Q,R)$.

\emph{Left side.} On $\partial S$ we have $z=f(x,y)$, hence $\dd z=f_x\dd x+f_y\dd y$, so
\[
\oint_{\partial S}P\,\dd x+Q\,\dd y+R\,\dd z=\oint_{\partial D}\bigl(P+Rf_x\bigr)\dd x+\bigl(Q+Rf_y\bigr)\dd y,
\]
where $P,Q,R$ are evaluated at $(x,y,f(x,y))$. By Green's theorem this equals
\[
\iint_D\Bigl[\pd{}x\bigl(Q+Rf_y\bigr)-\pd{}y\bigl(P+Rf_x\bigr)\Bigr]\dd A .
\]
\emph{Compute the integrand} with the chain rule, noting that $P,Q,R$ depend on $x$ through $x$ itself and through $z=f(x,y)$:
\begin{align*}
\pd{}x(Q+Rf_y)&=Q_x+Q_zf_x+(R_x+R_zf_x)f_y+Rf_{yx},\\
\pd{}y(P+Rf_x)&=P_y+P_zf_y+(R_y+R_zf_y)f_x+Rf_{xy}.
\end{align*}
Subtract. The terms $R_zf_xf_y$ cancel, and the terms $Rf_{yx}$ and $Rf_{xy}$ cancel by Clairaut's theorem (as $f\in C^2$). What remains is
\[
(Q_x-P_y)+(Q_z-R_y)f_x+(R_x-P_z)f_y .
\]
\emph{Right side.} The upward normal vector is $\vec N=(-f_x,-f_y,1)$ and $\hat{\vec n}\,\dd S=\vec N\,\dd A$ (Theorem~\ref{thm:fluxgraph}). With $\curl\vec F=(R_y-Q_z,\ P_z-R_x,\ Q_x-P_y)$,
\[
\curl\vec F\cdot\vec N=-(R_y-Q_z)f_x-(P_z-R_x)f_y+(Q_x-P_y),
\]
which is the same expression. Hence the two sides agree.
\end{proof}

A general surface is cut into pieces which are graphs over one of the three coordinate planes (or parametrized patches, using the same computation with $\vec r(u,v)$ in place of the graph). On a common edge of two pieces the boundary integrals cancel, because the edge is traversed in opposite directions. Hence Stokes' theorem holds for any surface that can be cut into such pieces.

\begin{example}[Using Stokes to replace a surface]
Let $\vec F=(y,\ z,\ x)$ and let $S$ be the part of the paraboloid $z=1-x^2-y^2$ with $z\ge0$, oriented upward. Then $\curl\vec F=(-1,-1,-1)$. Instead of integrating over the paraboloid, use $\partial S$: the unit circle $\vec r=(\cos t,\sin t,0)$. Here $\vec F\cdot\vec r\,'=(\sin t,0,\cos t)\cdot(-\sin t,\cos t,0)=-\sin^2t$, so
\[
\oint_{\partial S}\vec F\cdot\dd\vec r=-\int_0^{2\pi}\sin^2t\,\dd t=-\pi .
\]
\end{example}

\begin{theorem}[Consequences]\label{thm:consequences}
\begin{enumerate}[label=(\alph*)]
  \item If $\curl\vec F=\vec0$ on a simply connected region of $\R^3$ then $\vec F$ is conservative.
  \item If two surfaces $S_1,S_2$ have the same oriented boundary curve, then $\iint_{S_1}\curl\vec F\cdot\hat{\vec n}\dd S=\iint_{S_2}\curl\vec F\cdot\hat{\vec n}\dd S$. In particular, the flux of a curl through a \emph{closed} surface is $0$.
\end{enumerate}
\end{theorem}

\begin{proof}
(b) Both integrals equal $\oint_{\partial S}\vec F\cdot\dd\vec r$. For a closed surface, apply this to the two halves, or use the divergence theorem together with $\divg\curl\vec F=0$: $\oiint\curl\vec F\cdot\hat{\vec n}\dd S=\iiint\divg\curl\vec F\,\dd V=0$. (a) Let $C$ be a closed curve. In a simply connected region it is the boundary of some surface $S$ (this is what ``simply connected'' means in $\R^3$, informally), so by Stokes $\oint_C\vec F\cdot\dd\vec r=\iint_S\curl\vec F\cdot\hat{\vec n}\,\dd S=0$. By Theorem~\ref{thm:pathind}, $\vec F$ is conservative.
\end{proof}

\begin{theorem}[Curl as circulation density]
$\curl\vec F(\vec p)\cdot\hat{\vec n}=\lim_{\varepsilon\to0}\dfrac1{\pi\varepsilon^2}\oint_{C_\varepsilon}\vec F\cdot\dd\vec r$, where $C_\varepsilon$ is the circle of radius $\varepsilon$ around $\vec p$ in the plane with normal $\hat{\vec n}$.
\end{theorem}

\begin{proof}
By Stokes, the right side is the average of $\curl\vec F\cdot\hat{\vec n}$ over the disc, which tends to its value at $\vec p$ by continuity.
\end{proof}

%====================================================================
\section{The unifying picture}
%====================================================================

\begin{center}
\renewcommand{\arraystretch}{1.4}
\begin{tabular}{llll}
\toprule
\textbf{Theorem} & \textbf{Region} & \textbf{Derivative inside} & \textbf{Boundary} \\
\midrule
Fundamental theorem of calculus & interval $[a,b]$ & $f'$ & endpoints $\{a,b\}$ \\
Fundamental theorem for line integrals & curve $C$ & $\nabla f\cdot\dd\vec r$ & endpoints \\
Green & plane region $D$ & $Q_x-P_y$ & curve $\partial D$ \\
Stokes & surface $S$ & $\curl\vec F$ & curve $\partial S$ \\
Gauss & solid $E$ & $\divg\vec F$ & surface $\partial E$ \\
\bottomrule
\end{tabular}
\end{center}

All five are special cases of one statement, the generalized Stokes theorem $\int_\Omega\dd\omega=\int_{\partial\Omega}\omega$, which we prove in Chapter~\ref{ch:forms}. Two features are common to all of them: (1) \emph{the boundary of a region has no boundary} (the boundary of a disc is a circle, which has no endpoints), and this fact is the reason behind $\curl\nabla f=0$ and $\divg\curl\vec F=0$; (2) \emph{orientation} matters: the sign convention of the boundary is exactly what makes the interior cancellations work.

\begin{figure}[ht]
\centering
\begin{tikzpicture}[scale=1.0]
  \tikzset{blk/.style={draw=accent,fill=accent!10,rounded corners=2pt,minimum width=3.2cm,minimum height=0.9cm,align=center,font=\small}}
  \node[blk] (a) at (0,0) {$f$};
  \node[blk] (b) at (4.2,0) {$\nabla f$};
  \node[blk] (c) at (8.4,0) {$\curl\nabla f=0$};
  \draw[->,thick] (a)--(b) node[midway,above,font=\scriptsize]{grad};
  \draw[->,thick] (b)--(c) node[midway,above,font=\scriptsize]{curl};
  \node[blk,fill=accent2!12,draw=accent2] (d) at (0,-1.8) {$\vec A$};
  \node[blk,fill=accent2!12,draw=accent2] (e) at (4.2,-1.8) {$\curl\vec A$};
  \node[blk,fill=accent2!12,draw=accent2] (f) at (8.4,-1.8) {$\divg\curl\vec A=0$};
  \draw[->,thick] (d)--(e) node[midway,above,font=\scriptsize]{curl};
  \draw[->,thick] (e)--(f) node[midway,above,font=\scriptsize]{div};
\end{tikzpicture}
\caption{Applying two derivatives in a row gives zero. This is the algebraic shadow of ``the boundary of a boundary is empty''.}
\label{fig:dd0}
\end{figure}

%====================================================================
\section{Green's identities}
%====================================================================

The divergence theorem applied to a product gives the integration by parts formulas in several variables. They are the main tools for partial differential equations.

\begin{theorem}[Integration by parts]\label{thm:ibp}
For a $C^1$ scalar function $u$ and a $C^1$ vector field $\vec F$ on $E\cup\partial E$,
\[
\iiint_Eu\,\divg\vec F\,\dd V=\oiint_{\partial E}u\,\vec F\cdot\hat{\vec n}\,\dd S-\iiint_E\nabla u\cdot\vec F\,\dd V .
\]
\end{theorem}

\begin{proof}
By the product rule, $\divg(u\vec F)=\nabla u\cdot\vec F+u\divg\vec F$ (Theorem~\ref{thm:identities}). Integrate over $E$ and apply the divergence theorem to $u\vec F$.
\end{proof}

\begin{theorem}[Green's identities]\label{thm:greenid}
Let $u,v$ be $C^2$ on $E\cup\partial E$. Then
\begin{align}
\iiint_E\bigl(u\lap v+\nabla u\cdot\nabla v\bigr)\dd V&=\oiint_{\partial E}u\,\pd vn\,\dd S&&\text{(first identity)},\label{eq:green1}\\
\iiint_E\bigl(u\lap v-v\lap u\bigr)\dd V&=\oiint_{\partial E}\Bigl(u\pd vn-v\pd un\Bigr)\dd S&&\text{(second identity)},\label{eq:green2}
\end{align}
where $\pd vn=\nabla v\cdot\hat{\vec n}$ is the outward normal derivative.
\end{theorem}

\begin{proof}
For \eqref{eq:green1}, apply the divergence theorem to $\vec F=u\nabla v$; here $\divg(u\nabla v)=\nabla u\cdot\nabla v+u\lap v$. For \eqref{eq:green2}, write \eqref{eq:green1} again with $u$ and $v$ exchanged and subtract the two equations: the symmetric terms $\nabla u\cdot\nabla v$ cancel.
\end{proof}

\begin{corollary}[Energy and uniqueness]\label{cor:uniq}
Let $u$ be $C^2$ with $\lap u=0$ in $E$ and either $u=0$ or $\partial u/\partial n=0$ on all of $\partial E$. Then $u$ is constant in $E$ (and $u\equiv0$ in the first case if $E$ is connected).
\end{corollary}

\begin{proof}
Take $v=u$ in \eqref{eq:green1}: $\iiint_E\norm{\nabla u}^2\dd V=\oiint u\,\partial_nu\,\dd S=0$. A continuous nonnegative integrand with zero integral is zero, so $\nabla u\equiv\vec0$ and $u$ is constant on the connected set $E$; for $u=0$ on the boundary the constant is $0$.
\end{proof}

\begin{theorem}[Uniqueness for the Poisson equation]\label{thm:poissonuniq}
Let $E$ be bounded and connected. If $u_1,u_2$ are $C^2$ solutions of $\lap u=\rho$ in $E$ with the same values on $\partial E$ (Dirichlet condition), then $u_1=u_2$.
\end{theorem}

\begin{proof}
$w=u_1-u_2$ satisfies $\lap w=0$ and $w=0$ on $\partial E$. Apply Corollary~\ref{cor:uniq}.
\end{proof}

\begin{intuition}
Given the charge density $\rho$, the potential is completely determined by its values on the boundary. There is no room for another solution; this is why electrostatics problems with conductors at fixed potentials have a unique answer.
\end{intuition}

%====================================================================
\section{More consequences}
%====================================================================

\begin{theorem}[Reynolds-type identity: the surface integral of the gradient]\label{thm:gradint}
For a $C^1$ scalar function $f$, $\iiint_E\nabla f\,\dd V=\oiint_{\partial E}f\,\hat{\vec n}\,\dd S$, and $\iiint_E\curl\vec F\,\dd V=\oiint_{\partial E}\hat{\vec n}\times\vec F\,\dd S$.
\end{theorem}

\begin{proof}
Apply the divergence theorem to $\vec F=f\vec c$ for a constant vector $\vec c$: $\divg(f\vec c)=\nabla f\cdot\vec c$ and $f\vec c\cdot\hat{\vec n}=\vec c\cdot f\hat{\vec n}$, so $\vec c\cdot\bigl[\iiint\nabla f\dd V-\oiint f\hat{\vec n}\dd S\bigr]=0$ for every $\vec c$, which proves the first identity. For the second, apply the divergence theorem to $\vec F\times\vec c$: $\divg(\vec F\times\vec c)=\vec c\cdot\curl\vec F$ (identity 6 of Theorem~\ref{thm:identities}, with $\curl\vec c=0$), and $(\vec F\times\vec c)\cdot\hat{\vec n}=\vec c\cdot(\hat{\vec n}\times\vec F)$ (cyclic property of the triple product). Again $\vec c$ is arbitrary.
\end{proof}

\begin{example}[Archimedes' principle]
A body $E$ is immersed in a fluid of density $\rho_0$; the pressure at depth is $p=-\rho_0gz$ ($z$ up, $p=0$ at $z=0$). The force of the fluid on the body is $-\oiint p\,\hat{\vec n}\,\dd S$. By Theorem~\ref{thm:gradint}, this is $-\iiint\nabla p\,\dd V=\iiint\rho_0g\,\khat\,\dd V=\rho_0g\Vol(E)\,\khat$, which is the weight of the displaced fluid, pointing up.
\end{example}

%====================================================================
\begin{summarybox}
\begin{itemize}
  \item \textbf{Green:} $\oint_{\partial D}P\dd x+Q\dd y=\iint_D(Q_x-P_y)\dd A$; the region lies on the left.
  \item \textbf{Divergence theorem:} $\iiint_E\divg\vec F\dd V=\oiint_{\partial E}\vec F\cdot\hat{\vec n}\dd S$ with outward $\hat{\vec n}$.
  \item \textbf{Stokes:} $\oint_{\partial S}\vec F\cdot\dd\vec r=\iint_S\curl\vec F\cdot\hat{\vec n}\dd S$, orientation by the right-hand rule.
  \item Proofs: reduce to simple regions, use the fundamental theorem of calculus in one variable and Fubini; glue pieces together, where interior boundaries cancel because of opposite orientations.
  \item Consequences: divergence and curl as limits of flux and circulation densities; Gauss's law ($4\pi$); Green's identities; uniqueness for $\lap u=\rho$; Archimedes.
  \item The singular field $\vec r/r^3$ or the vortex is handled by cutting out a small ball or disc.
\end{itemize}
\end{summarybox}

\begin{interview}
\begin{enumerate}
  \item \emph{State Green's, Stokes' and Gauss' theorems and how they are related.} Each is $\int_\Omega\dd\omega=\int_{\partial\Omega}\omega$; Green is Stokes in a plane.
  \item \emph{Prove the divergence theorem for a simple region.} Reduce to $\iiint R_z=\oiint Rn_z$, use the fundamental theorem of calculus in $z$ and compute the top, bottom and side contributions.
  \item \emph{Why can interior boundaries be ignored?} They appear twice with opposite orientation (outward normals are opposite) and cancel.
  \item \emph{Compute the flux of $\vec r/r^3$ through a closed surface.} $4\pi$ if the origin is inside, $0$ otherwise.
  \item \emph{Why is the flux of a curl through a closed surface zero?} $\divg\curl=0$ and the divergence theorem.
  \item \emph{Prove uniqueness of the Poisson problem with Dirichlet data.} Green's first identity with $v=u=w$.
  \item \emph{Orientation conventions.} Region on the left for Green; right-hand rule for Stokes; outward normal for Gauss.
\end{enumerate}
\end{interview}

\subsection*{Exercises}
\begin{exercise}
Use Green's theorem to compute $\oint_C(x^2y\,\dd x+xy^2\,\dd y)$ for the boundary $C$ of the unit disc. (Answer: $\pi/2$.)
\end{exercise}
\begin{exercise}
Compute the outward flux of $\vec F=(xy,\ yz,\ zx)$ through the boundary of the unit cube $[0,1]^3$. (Answer: $\tfrac32$.)
\end{exercise}
\begin{exercise}
Verify Stokes' theorem for $\vec F=(-y,x,z)$ and the upper unit hemisphere with its boundary circle. (Both sides: $2\pi$.)
\end{exercise}
\begin{exercise}
Show that $\oiint_{\partial E}\hat{\vec n}\,\dd S=\vec0$ for any closed surface. (Hint: Theorem~\ref{thm:gradint} with $f=1$.)
\end{exercise}
\begin{exercise}
If $\lap u=0$ in $E$, show $\oiint_{\partial E}\partial_nu\,\dd S=0$.
\end{exercise}
\begin{exercise}
Let $\vec F=\dfrac{(-y,x)}{x^2+y^2}$ and $C$ any simple closed curve around the origin, counterclockwise. Use Theorem~\ref{thm:greensing} to show $\oint_C\vec F\cdot\dd\vec r=2\pi$.
\end{exercise}
\begin{exercise}
(Green's second identity in 2D with $v=\ln r$.) Derive the mean value property: for $u$ harmonic in a disc, $u(\vec 0)$ is the average of $u$ over any circle $r=\rho$ in the disc. (Hint: apply the identity on the annulus $\varepsilon\le r\le\rho$ with $v=\ln r$ and let $\varepsilon\to0$.)
\end{exercise}
